\chapter*{第四部分收束：从经典案例走向深度学习}
\addcontentsline{toc}{chapter}{第四部分收束：从经典案例走向深度学习}

Part IV 的四个案例展示了一个共同事实：很多天文与物理问题并不是一开始就需要深度学习。只要数据规模适中、特征含义清楚、baseline 合理，传统机器学习和清楚的诊断图表就已经能完成大量有价值的工作。

但这些案例也自然暴露出深度学习的入口。星系形态图像中，手工特征很难完全表达旋臂、棒结构和局部纹理；光谱分类中，局部谱线模式可能比全局手工特征更有信息；时间序列和序列化光谱中，局部模式、长程依赖和表示学习会越来越重要。

\section*{什么时候需要更复杂模型}

更复杂模型通常不是因为“更先进”，而是因为简单表示开始吃力：

\begin{itemize}
    \item 图像中有局部结构，手工特征难以覆盖；
    \item 光谱中谱线位置会轻微移动，raw-feature baseline 不够稳；
    \item 样本结构复杂，线性或树模型只能抓到局部规则；
    \item 你希望模型自己学习中间表示，而不是完全依赖人工特征；
    \item 任务需要处理序列、patch、token 或多模态输入。
\end{itemize}

同时，复杂模型也会带来新的责任：更多参数、更高过拟合风险、更难解释、更强算力依赖，以及更严格的数据划分要求。

\section*{进入 Part V 前的基线意识}

Part V 会进入神经网络、卷积、自编码器、Transformer 和现代 AI 工具。但请把 Part IV 的经验带过去：\textbf{深度学习不应该抹掉 baseline，而应该站在 baseline 之上。}

一个健康的深度学习实验，至少应能回答：

\begin{itemize}
    \item 简单模型已经做到什么程度；
    \item 深度模型具体解决了哪个表示问题；
    \item 改进是否来自真正的结构学习，而不是数据泄漏或调参偶然；
    \item 错误样本是否仍然能被人类复查；
    \item 模型输出是否能回到科学问题中解释。
\end{itemize}

如果这些问题答不上来，复杂模型就只是更漂亮的黑箱。
